\section{The \texttt{Mirror\_curved} McXtrace Component}
A cylindrically curved mirror (in YZ)

\subsection*{Identification}
\begin{itemize}
  \item \textbf{Author:} Erik Knudsen
  \item \textbf{Origin:} Risoe
  \item \textbf{Date:} September 25th, 2009
\end{itemize}

\subsection*{Description}
mirror is in the YZ-plane curved towards positive X if radius is positive

Example: Mirror\_curved( radius=2, length=20e-3, width=40e-3, coating="AlMgF2\_disco.dat")

\subsection*{Input parameters}
Parameters in \textbf{boldface} are required; the others are optional.

\begin{longtable}{p{0.22\textwidth}p{0.12\textwidth}p{0.46\textwidth}p{0.14\textwidth}}
\toprule
\textbf{Name} & \textbf{Unit} & \textbf{Description} & \textbf{Default} \\
\midrule
\endhead
radius & m & Radius of curvature, along Z. & 1 \\
R0 & 1 & Constant reflectivity value (mostly relevant for debugging, when coating=""). & 1 \\
coating & str & Name of file containing the material data (i.e. f1 and f2) for the coating & "Be.txt" \\
zdepth & m & Length of the unbent mirror along Z. & 0.2 \\
yheight & m & Width of the mirror along Y. & 0.2 \\
length & m & Length of the unbent mirror along Z = zdepth & 0 \\
width & m & Width of the mirror along Y = yheight & 0 \\
\bottomrule
\end{longtable}

\subsection*{Links}
\begin{itemize}
  \item Component source code found in file \texttt{Mirror\_curved.comp}.
\end{itemize}
\IfFileExists{optics/Mirror_curved_static.tex}{\input{optics/Mirror_curved_static.tex}}{}